\documentclass[11pt,a4paper]{article}
% main (standalone preamble for editing)
% vim: nowrap: spell:
% ══════════════════════════════════════════════════════════════════════════════
% Speed Maths --- shared preamble
% Included by every sheet and answer file via:  % main (standalone preamble for editing)
% vim: nowrap: spell:
% ══════════════════════════════════════════════════════════════════════════════
% Speed Maths --- shared preamble
% Included by every sheet and answer file via:  % main (standalone preamble for editing)
% vim: nowrap: spell:
% ══════════════════════════════════════════════════════════════════════════════
% Speed Maths --- shared preamble
% Included by every sheet and answer file via:  \input{../../shared/preamble}
% Editing THIS file changes the look of every sheet/answer across every pillar.
% ══════════════════════════════════════════════════════════════════════════════

\usepackage[a4paper, top=25mm, bottom=25mm, left=22mm, right=22mm]{geometry}
\usepackage{amsmath, amssymb}
\usepackage{enumitem}
\usepackage{booktabs}
\usepackage{array}
\usepackage{parskip}
\usepackage{microtype}
\usepackage{fancyhdr}
\usepackage{titlesec}
\usepackage{mdframed}
\usepackage{xcolor}
\usepackage[hidelinks]{hyperref}
\usepackage{graphicx}
\usepackage{tikz}
\usepackage{pgfplots}
\pgfplotsset{compat=1.18}

% ── Colours ───────────────────────────────────────────────────────────────────
\definecolor{darkgray}{gray}{0.15}
\definecolor{medgray}{gray}{0.45}
\definecolor{lightgray}{gray}{0.93}
\definecolor{boxgray}{gray}{0.88}

% ── Section formatting ──────────────────────────────────────────────────────────
\titleformat{\section}[block]
  {\normalfont\large\bfseries}{}
  {0pt}{}[\vspace{2pt}\hrule\vspace{6pt}]
\titlespacing*{\section}{0pt}{14pt}{4pt}

% ── List formatting ─────────────────────────────────────────────────────────────
\setlist[enumerate,1]{label=\textbf{\arabic*.}, leftmargin=2em, itemsep=10pt}

% ── Branding constants (edit here, not per-file) ────────────────────────────────
\newcommand{\SpeedExamLine}{\textbf{Speed Maths:} Competition \& Exam Prep}
\newcommand{\SpeedCredit}{Series by \href{https://www.linkedin.com/in/cerealdev/}{David Akinyele-Aje}}

% ── Header / footer ──────────────────────────────────────────────────────────────
% Usage (once, after \input{../../shared/preamble} and before \begin{document}):
%   \SpeedHeader{Algebra}{3}
% First arg  = pillar name shown as "Speed <Pillar> --- Daily Drill"
% Second arg = worksheet number
\pagestyle{fancy}
\newcommand{\SpeedHeader}[2]{%
  \fancyhf{}%
  \renewcommand{\headrulewidth}{0.4pt}%
  \setlength{\headheight}{14pt}%
  \fancyhead[L]{\small\textbf{Speed #1 --- Daily Drill}}%
  \fancyhead[R]{\small Worksheet \##2}%
  \fancyfoot[L]{\small \SpeedExamLine}%
  \fancyfoot[C]{\small\thepage}%
  \fancyfoot[R]{\small \SpeedCredit}%
  \renewcommand{\footrulewidth}{0.4pt}%
}

% ── Title block (used at the top of every sheet AND answer file) ───────────────
% Usage: \SpeedTitleBlock{Daily Algebra Drill \#3}{Jane Doe}
% %% AGENTS: When generating a new sheet, ask the human contributor for the name they want credited in argument 2.
% For answer files, pass the "--- Answers \& Investigations" suffix in the arg.
\newcommand{\SpeedTitleBlock}[2]{%
  \begin{center}
    {\LARGE\bfseries #1}\\[4pt]
    {\normalsize\itshape Sheet by #2}\\[6pt]
    \hrule\vspace{4pt}
    \begin{tabular}{llcc}
      \toprule
      \textbf{Section} & \textbf{Topic} & \textbf{Questions} & \textbf{Target time}\\
      \midrule
      A & Rapid Recognition          & 10 & 2 min 30 sec \\
      B & Manipulation Drills        & 10 & 8 min        \\
      C & Substitution \& Structure  & 8  & 10 min       \\
      D & Challenge                  & 5  & 15 min       \\
      \bottomrule
    \end{tabular}
    \vspace{4pt}
    \hrule
  \end{center}
}

% ── Sheet metadata: topic + the day's new tools ────────────────────────────────
% Usage (immediately after \SpeedTitleBlock, sheets only):
%   \SpeedMeta{Expanding, factorising, and the identity toolkit}
%             {difference of squares, sum/product form, completing the square}
%
% Arg 1 = the sheet's topic in one line. The website lists this instead of
%         "Sheet 03", so write the thing a reader would actually search for.
% Arg 2 = comma-separated, at most five: the tools this sheet introduces that
%         earlier sheets in the pillar did not. These become the website's tags.
%         Leave out anything the reader already met — the tags are meant to say
%         what is new today, not to summarise the sheet.
%
% Both are parsed straight out of this file by tools/build_website.py, so the
% sheet stays the single source of truth and the site cannot drift from it.
%
% This renders NOTHING. It is a declaration for the build script, not a piece of
% the sheet's layout: adding it to a sheet must not change that sheet's PDF, so
% the 25 existing sheets can be annotated without a single one of them being
% reissued. If the toolkit should also appear on the page, that is a separate
% decision about the sheet design — make it deliberately, here, not by leaving
% this macro printing something nobody asked it to.
\newcommand{\SpeedMeta}[2]{}

% ── Answer / Method / Investigate-further boxes (answer files only) ────────────
\newmdenv[
  backgroundcolor=lightgray,
  linecolor=medgray,
  linewidth=0.5pt,
  leftmargin=0pt, rightmargin=0pt,
  innerleftmargin=8pt, innerrightmargin=8pt,
  innertopmargin=5pt, innerbottommargin=5pt,
  skipabove=4pt, skipbelow=4pt
]{ansbox}

\newmdenv[
  backgroundcolor=white,
  linecolor=darkgray,
  linewidth=0.8pt,
  leftline=true, rightline=false, topline=false, bottomline=false,
  leftmargin=0pt, rightmargin=0pt,
  innerleftmargin=8pt, innerrightmargin=6pt,
  innertopmargin=4pt, innerbottommargin=4pt,
  skipabove=4pt, skipbelow=6pt
]{invbox}

\newcommand{\ans}[1]{%
  \begin{ansbox}
    \textbf{Answer:} #1
  \end{ansbox}}

\newcommand{\method}[1]{%
  {\small\textbf{Method:} #1}\par}

\newcommand{\inv}[1]{%
  \begin{invbox}
    \textbf{\small Investigate further:} {\small #1}
  \end{invbox}}

% ── Misc helpers ────────────────────────────────────────────────────────────────
\newcommand{\qn}[1]{\textbf{#1.}\;}

% Mistake-linking: attach a Top 5 Patterns / Common Traps bullet to the question
% label(s) in this sheet where it actually shows up, e.g. \seealso{B6, D2}
\newcommand{\seealso}[1]{\hfill\textit{\footnotesize(see #1)}}

% ── Graph options macro ─────────────────────────────────────────────────────────
% Usage: \GraphOptions{tikz code for A}{tikz code for B}{tikz code for C}{tikz code for D}
\newcommand{\GraphOptions}[4]{%
  \begin{center}
    \begin{tabular}{cc}
      \textbf{A)} & \textbf{B)} \\
      #1 & #2 \\[10pt]
      \textbf{C)} & \textbf{D)} \\
      #3 & #4
    \end{tabular}
  \end{center}
}

% ── Closing motivational rule + cross-promo footer (sheets only) ────────────────
% Usage: \SpeedClosing{"Quote text."}
\newcommand{\SpeedClosing}[1]{%
  \vspace{20pt}
  \begin{center}
  \rule{0.6\textwidth}{0.4pt}\\[4pt]
  {\small\itshape #1}\\[4pt]
  \rule{0.6\textwidth}{0.4pt}
  \end{center}
  \begin{center}
  {\small Found a mistake or want to contribute a sheet?}\\
  {\small Visit the open-source project at \href{https://github.com/The-CerealDev/speed-maths}{github.com/The-CerealDev/speed-maths}}
  \end{center}
}

% Editing THIS file changes the look of every sheet/answer across every pillar.
% ══════════════════════════════════════════════════════════════════════════════

\usepackage[a4paper, top=25mm, bottom=25mm, left=22mm, right=22mm]{geometry}
\usepackage{amsmath, amssymb}
\usepackage{enumitem}
\usepackage{booktabs}
\usepackage{array}
\usepackage{parskip}
\usepackage{microtype}
\usepackage{fancyhdr}
\usepackage{titlesec}
\usepackage{mdframed}
\usepackage{xcolor}
\usepackage[hidelinks]{hyperref}
\usepackage{graphicx}
\usepackage{tikz}
\usepackage{pgfplots}
\pgfplotsset{compat=1.18}

% ── Colours ───────────────────────────────────────────────────────────────────
\definecolor{darkgray}{gray}{0.15}
\definecolor{medgray}{gray}{0.45}
\definecolor{lightgray}{gray}{0.93}
\definecolor{boxgray}{gray}{0.88}

% ── Section formatting ──────────────────────────────────────────────────────────
\titleformat{\section}[block]
  {\normalfont\large\bfseries}{}
  {0pt}{}[\vspace{2pt}\hrule\vspace{6pt}]
\titlespacing*{\section}{0pt}{14pt}{4pt}

% ── List formatting ─────────────────────────────────────────────────────────────
\setlist[enumerate,1]{label=\textbf{\arabic*.}, leftmargin=2em, itemsep=10pt}

% ── Branding constants (edit here, not per-file) ────────────────────────────────
\newcommand{\SpeedExamLine}{\textbf{Speed Maths:} Competition \& Exam Prep}
\newcommand{\SpeedCredit}{Series by \href{https://www.linkedin.com/in/cerealdev/}{David Akinyele-Aje}}

% ── Header / footer ──────────────────────────────────────────────────────────────
% Usage (once, after % main (standalone preamble for editing)
% vim: nowrap: spell:
% ══════════════════════════════════════════════════════════════════════════════
% Speed Maths --- shared preamble
% Included by every sheet and answer file via:  \input{../../shared/preamble}
% Editing THIS file changes the look of every sheet/answer across every pillar.
% ══════════════════════════════════════════════════════════════════════════════

\usepackage[a4paper, top=25mm, bottom=25mm, left=22mm, right=22mm]{geometry}
\usepackage{amsmath, amssymb}
\usepackage{enumitem}
\usepackage{booktabs}
\usepackage{array}
\usepackage{parskip}
\usepackage{microtype}
\usepackage{fancyhdr}
\usepackage{titlesec}
\usepackage{mdframed}
\usepackage{xcolor}
\usepackage[hidelinks]{hyperref}
\usepackage{graphicx}
\usepackage{tikz}
\usepackage{pgfplots}
\pgfplotsset{compat=1.18}

% ── Colours ───────────────────────────────────────────────────────────────────
\definecolor{darkgray}{gray}{0.15}
\definecolor{medgray}{gray}{0.45}
\definecolor{lightgray}{gray}{0.93}
\definecolor{boxgray}{gray}{0.88}

% ── Section formatting ──────────────────────────────────────────────────────────
\titleformat{\section}[block]
  {\normalfont\large\bfseries}{}
  {0pt}{}[\vspace{2pt}\hrule\vspace{6pt}]
\titlespacing*{\section}{0pt}{14pt}{4pt}

% ── List formatting ─────────────────────────────────────────────────────────────
\setlist[enumerate,1]{label=\textbf{\arabic*.}, leftmargin=2em, itemsep=10pt}

% ── Branding constants (edit here, not per-file) ────────────────────────────────
\newcommand{\SpeedExamLine}{\textbf{Speed Maths:} Competition \& Exam Prep}
\newcommand{\SpeedCredit}{Series by \href{https://www.linkedin.com/in/cerealdev/}{David Akinyele-Aje}}

% ── Header / footer ──────────────────────────────────────────────────────────────
% Usage (once, after \input{../../shared/preamble} and before \begin{document}):
%   \SpeedHeader{Algebra}{3}
% First arg  = pillar name shown as "Speed <Pillar> --- Daily Drill"
% Second arg = worksheet number
\pagestyle{fancy}
\newcommand{\SpeedHeader}[2]{%
  \fancyhf{}%
  \renewcommand{\headrulewidth}{0.4pt}%
  \setlength{\headheight}{14pt}%
  \fancyhead[L]{\small\textbf{Speed #1 --- Daily Drill}}%
  \fancyhead[R]{\small Worksheet \##2}%
  \fancyfoot[L]{\small \SpeedExamLine}%
  \fancyfoot[C]{\small\thepage}%
  \fancyfoot[R]{\small \SpeedCredit}%
  \renewcommand{\footrulewidth}{0.4pt}%
}

% ── Title block (used at the top of every sheet AND answer file) ───────────────
% Usage: \SpeedTitleBlock{Daily Algebra Drill \#3}{Jane Doe}
% %% AGENTS: When generating a new sheet, ask the human contributor for the name they want credited in argument 2.
% For answer files, pass the "--- Answers \& Investigations" suffix in the arg.
\newcommand{\SpeedTitleBlock}[2]{%
  \begin{center}
    {\LARGE\bfseries #1}\\[4pt]
    {\normalsize\itshape Sheet by #2}\\[6pt]
    \hrule\vspace{4pt}
    \begin{tabular}{llcc}
      \toprule
      \textbf{Section} & \textbf{Topic} & \textbf{Questions} & \textbf{Target time}\\
      \midrule
      A & Rapid Recognition          & 10 & 2 min 30 sec \\
      B & Manipulation Drills        & 10 & 8 min        \\
      C & Substitution \& Structure  & 8  & 10 min       \\
      D & Challenge                  & 5  & 15 min       \\
      \bottomrule
    \end{tabular}
    \vspace{4pt}
    \hrule
  \end{center}
}

% ── Sheet metadata: topic + the day's new tools ────────────────────────────────
% Usage (immediately after \SpeedTitleBlock, sheets only):
%   \SpeedMeta{Expanding, factorising, and the identity toolkit}
%             {difference of squares, sum/product form, completing the square}
%
% Arg 1 = the sheet's topic in one line. The website lists this instead of
%         "Sheet 03", so write the thing a reader would actually search for.
% Arg 2 = comma-separated, at most five: the tools this sheet introduces that
%         earlier sheets in the pillar did not. These become the website's tags.
%         Leave out anything the reader already met — the tags are meant to say
%         what is new today, not to summarise the sheet.
%
% Both are parsed straight out of this file by tools/build_website.py, so the
% sheet stays the single source of truth and the site cannot drift from it.
%
% This renders NOTHING. It is a declaration for the build script, not a piece of
% the sheet's layout: adding it to a sheet must not change that sheet's PDF, so
% the 25 existing sheets can be annotated without a single one of them being
% reissued. If the toolkit should also appear on the page, that is a separate
% decision about the sheet design — make it deliberately, here, not by leaving
% this macro printing something nobody asked it to.
\newcommand{\SpeedMeta}[2]{}

% ── Answer / Method / Investigate-further boxes (answer files only) ────────────
\newmdenv[
  backgroundcolor=lightgray,
  linecolor=medgray,
  linewidth=0.5pt,
  leftmargin=0pt, rightmargin=0pt,
  innerleftmargin=8pt, innerrightmargin=8pt,
  innertopmargin=5pt, innerbottommargin=5pt,
  skipabove=4pt, skipbelow=4pt
]{ansbox}

\newmdenv[
  backgroundcolor=white,
  linecolor=darkgray,
  linewidth=0.8pt,
  leftline=true, rightline=false, topline=false, bottomline=false,
  leftmargin=0pt, rightmargin=0pt,
  innerleftmargin=8pt, innerrightmargin=6pt,
  innertopmargin=4pt, innerbottommargin=4pt,
  skipabove=4pt, skipbelow=6pt
]{invbox}

\newcommand{\ans}[1]{%
  \begin{ansbox}
    \textbf{Answer:} #1
  \end{ansbox}}

\newcommand{\method}[1]{%
  {\small\textbf{Method:} #1}\par}

\newcommand{\inv}[1]{%
  \begin{invbox}
    \textbf{\small Investigate further:} {\small #1}
  \end{invbox}}

% ── Misc helpers ────────────────────────────────────────────────────────────────
\newcommand{\qn}[1]{\textbf{#1.}\;}

% Mistake-linking: attach a Top 5 Patterns / Common Traps bullet to the question
% label(s) in this sheet where it actually shows up, e.g. \seealso{B6, D2}
\newcommand{\seealso}[1]{\hfill\textit{\footnotesize(see #1)}}

% ── Graph options macro ─────────────────────────────────────────────────────────
% Usage: \GraphOptions{tikz code for A}{tikz code for B}{tikz code for C}{tikz code for D}
\newcommand{\GraphOptions}[4]{%
  \begin{center}
    \begin{tabular}{cc}
      \textbf{A)} & \textbf{B)} \\
      #1 & #2 \\[10pt]
      \textbf{C)} & \textbf{D)} \\
      #3 & #4
    \end{tabular}
  \end{center}
}

% ── Closing motivational rule + cross-promo footer (sheets only) ────────────────
% Usage: \SpeedClosing{"Quote text."}
\newcommand{\SpeedClosing}[1]{%
  \vspace{20pt}
  \begin{center}
  \rule{0.6\textwidth}{0.4pt}\\[4pt]
  {\small\itshape #1}\\[4pt]
  \rule{0.6\textwidth}{0.4pt}
  \end{center}
  \begin{center}
  {\small Found a mistake or want to contribute a sheet?}\\
  {\small Visit the open-source project at \href{https://github.com/The-CerealDev/speed-maths}{github.com/The-CerealDev/speed-maths}}
  \end{center}
}
 and before \begin{document}):
%   \SpeedHeader{Algebra}{3}
% First arg  = pillar name shown as "Speed <Pillar> --- Daily Drill"
% Second arg = worksheet number
\pagestyle{fancy}
\newcommand{\SpeedHeader}[2]{%
  \fancyhf{}%
  \renewcommand{\headrulewidth}{0.4pt}%
  \setlength{\headheight}{14pt}%
  \fancyhead[L]{\small\textbf{Speed #1 --- Daily Drill}}%
  \fancyhead[R]{\small Worksheet \##2}%
  \fancyfoot[L]{\small \SpeedExamLine}%
  \fancyfoot[C]{\small\thepage}%
  \fancyfoot[R]{\small \SpeedCredit}%
  \renewcommand{\footrulewidth}{0.4pt}%
}

% ── Title block (used at the top of every sheet AND answer file) ───────────────
% Usage: \SpeedTitleBlock{Daily Algebra Drill \#3}{Jane Doe}
% %% AGENTS: When generating a new sheet, ask the human contributor for the name they want credited in argument 2.
% For answer files, pass the "--- Answers \& Investigations" suffix in the arg.
\newcommand{\SpeedTitleBlock}[2]{%
  \begin{center}
    {\LARGE\bfseries #1}\\[4pt]
    {\normalsize\itshape Sheet by #2}\\[6pt]
    \hrule\vspace{4pt}
    \begin{tabular}{llcc}
      \toprule
      \textbf{Section} & \textbf{Topic} & \textbf{Questions} & \textbf{Target time}\\
      \midrule
      A & Rapid Recognition          & 10 & 2 min 30 sec \\
      B & Manipulation Drills        & 10 & 8 min        \\
      C & Substitution \& Structure  & 8  & 10 min       \\
      D & Challenge                  & 5  & 15 min       \\
      \bottomrule
    \end{tabular}
    \vspace{4pt}
    \hrule
  \end{center}
}

% ── Sheet metadata: topic + the day's new tools ────────────────────────────────
% Usage (immediately after \SpeedTitleBlock, sheets only):
%   \SpeedMeta{Expanding, factorising, and the identity toolkit}
%             {difference of squares, sum/product form, completing the square}
%
% Arg 1 = the sheet's topic in one line. The website lists this instead of
%         "Sheet 03", so write the thing a reader would actually search for.
% Arg 2 = comma-separated, at most five: the tools this sheet introduces that
%         earlier sheets in the pillar did not. These become the website's tags.
%         Leave out anything the reader already met — the tags are meant to say
%         what is new today, not to summarise the sheet.
%
% Both are parsed straight out of this file by tools/build_website.py, so the
% sheet stays the single source of truth and the site cannot drift from it.
%
% This renders NOTHING. It is a declaration for the build script, not a piece of
% the sheet's layout: adding it to a sheet must not change that sheet's PDF, so
% the 25 existing sheets can be annotated without a single one of them being
% reissued. If the toolkit should also appear on the page, that is a separate
% decision about the sheet design — make it deliberately, here, not by leaving
% this macro printing something nobody asked it to.
\newcommand{\SpeedMeta}[2]{}

% ── Answer / Method / Investigate-further boxes (answer files only) ────────────
\newmdenv[
  backgroundcolor=lightgray,
  linecolor=medgray,
  linewidth=0.5pt,
  leftmargin=0pt, rightmargin=0pt,
  innerleftmargin=8pt, innerrightmargin=8pt,
  innertopmargin=5pt, innerbottommargin=5pt,
  skipabove=4pt, skipbelow=4pt
]{ansbox}

\newmdenv[
  backgroundcolor=white,
  linecolor=darkgray,
  linewidth=0.8pt,
  leftline=true, rightline=false, topline=false, bottomline=false,
  leftmargin=0pt, rightmargin=0pt,
  innerleftmargin=8pt, innerrightmargin=6pt,
  innertopmargin=4pt, innerbottommargin=4pt,
  skipabove=4pt, skipbelow=6pt
]{invbox}

\newcommand{\ans}[1]{%
  \begin{ansbox}
    \textbf{Answer:} #1
  \end{ansbox}}

\newcommand{\method}[1]{%
  {\small\textbf{Method:} #1}\par}

\newcommand{\inv}[1]{%
  \begin{invbox}
    \textbf{\small Investigate further:} {\small #1}
  \end{invbox}}

% ── Misc helpers ────────────────────────────────────────────────────────────────
\newcommand{\qn}[1]{\textbf{#1.}\;}

% Mistake-linking: attach a Top 5 Patterns / Common Traps bullet to the question
% label(s) in this sheet where it actually shows up, e.g. \seealso{B6, D2}
\newcommand{\seealso}[1]{\hfill\textit{\footnotesize(see #1)}}

% ── Graph options macro ─────────────────────────────────────────────────────────
% Usage: \GraphOptions{tikz code for A}{tikz code for B}{tikz code for C}{tikz code for D}
\newcommand{\GraphOptions}[4]{%
  \begin{center}
    \begin{tabular}{cc}
      \textbf{A)} & \textbf{B)} \\
      #1 & #2 \\[10pt]
      \textbf{C)} & \textbf{D)} \\
      #3 & #4
    \end{tabular}
  \end{center}
}

% ── Closing motivational rule + cross-promo footer (sheets only) ────────────────
% Usage: \SpeedClosing{"Quote text."}
\newcommand{\SpeedClosing}[1]{%
  \vspace{20pt}
  \begin{center}
  \rule{0.6\textwidth}{0.4pt}\\[4pt]
  {\small\itshape #1}\\[4pt]
  \rule{0.6\textwidth}{0.4pt}
  \end{center}
  \begin{center}
  {\small Found a mistake or want to contribute a sheet?}\\
  {\small Visit the open-source project at \href{https://github.com/The-CerealDev/speed-maths}{github.com/The-CerealDev/speed-maths}}
  \end{center}
}

% Editing THIS file changes the look of every sheet/answer across every pillar.
% ══════════════════════════════════════════════════════════════════════════════

\usepackage[a4paper, top=25mm, bottom=25mm, left=22mm, right=22mm]{geometry}
\usepackage{amsmath, amssymb}
\usepackage{enumitem}
\usepackage{booktabs}
\usepackage{array}
\usepackage{parskip}
\usepackage{microtype}
\usepackage{fancyhdr}
\usepackage{titlesec}
\usepackage{mdframed}
\usepackage{xcolor}
\usepackage[hidelinks]{hyperref}
\usepackage{graphicx}
\usepackage{tikz}
\usepackage{pgfplots}
\pgfplotsset{compat=1.18}

% ── Colours ───────────────────────────────────────────────────────────────────
\definecolor{darkgray}{gray}{0.15}
\definecolor{medgray}{gray}{0.45}
\definecolor{lightgray}{gray}{0.93}
\definecolor{boxgray}{gray}{0.88}

% ── Section formatting ──────────────────────────────────────────────────────────
\titleformat{\section}[block]
  {\normalfont\large\bfseries}{}
  {0pt}{}[\vspace{2pt}\hrule\vspace{6pt}]
\titlespacing*{\section}{0pt}{14pt}{4pt}

% ── List formatting ─────────────────────────────────────────────────────────────
\setlist[enumerate,1]{label=\textbf{\arabic*.}, leftmargin=2em, itemsep=10pt}

% ── Branding constants (edit here, not per-file) ────────────────────────────────
\newcommand{\SpeedExamLine}{\textbf{Speed Maths:} Competition \& Exam Prep}
\newcommand{\SpeedCredit}{Series by \href{https://www.linkedin.com/in/cerealdev/}{David Akinyele-Aje}}

% ── Header / footer ──────────────────────────────────────────────────────────────
% Usage (once, after % main (standalone preamble for editing)
% vim: nowrap: spell:
% ══════════════════════════════════════════════════════════════════════════════
% Speed Maths --- shared preamble
% Included by every sheet and answer file via:  % main (standalone preamble for editing)
% vim: nowrap: spell:
% ══════════════════════════════════════════════════════════════════════════════
% Speed Maths --- shared preamble
% Included by every sheet and answer file via:  \input{../../shared/preamble}
% Editing THIS file changes the look of every sheet/answer across every pillar.
% ══════════════════════════════════════════════════════════════════════════════

\usepackage[a4paper, top=25mm, bottom=25mm, left=22mm, right=22mm]{geometry}
\usepackage{amsmath, amssymb}
\usepackage{enumitem}
\usepackage{booktabs}
\usepackage{array}
\usepackage{parskip}
\usepackage{microtype}
\usepackage{fancyhdr}
\usepackage{titlesec}
\usepackage{mdframed}
\usepackage{xcolor}
\usepackage[hidelinks]{hyperref}
\usepackage{graphicx}
\usepackage{tikz}
\usepackage{pgfplots}
\pgfplotsset{compat=1.18}

% ── Colours ───────────────────────────────────────────────────────────────────
\definecolor{darkgray}{gray}{0.15}
\definecolor{medgray}{gray}{0.45}
\definecolor{lightgray}{gray}{0.93}
\definecolor{boxgray}{gray}{0.88}

% ── Section formatting ──────────────────────────────────────────────────────────
\titleformat{\section}[block]
  {\normalfont\large\bfseries}{}
  {0pt}{}[\vspace{2pt}\hrule\vspace{6pt}]
\titlespacing*{\section}{0pt}{14pt}{4pt}

% ── List formatting ─────────────────────────────────────────────────────────────
\setlist[enumerate,1]{label=\textbf{\arabic*.}, leftmargin=2em, itemsep=10pt}

% ── Branding constants (edit here, not per-file) ────────────────────────────────
\newcommand{\SpeedExamLine}{\textbf{Speed Maths:} Competition \& Exam Prep}
\newcommand{\SpeedCredit}{Series by \href{https://www.linkedin.com/in/cerealdev/}{David Akinyele-Aje}}

% ── Header / footer ──────────────────────────────────────────────────────────────
% Usage (once, after \input{../../shared/preamble} and before \begin{document}):
%   \SpeedHeader{Algebra}{3}
% First arg  = pillar name shown as "Speed <Pillar> --- Daily Drill"
% Second arg = worksheet number
\pagestyle{fancy}
\newcommand{\SpeedHeader}[2]{%
  \fancyhf{}%
  \renewcommand{\headrulewidth}{0.4pt}%
  \setlength{\headheight}{14pt}%
  \fancyhead[L]{\small\textbf{Speed #1 --- Daily Drill}}%
  \fancyhead[R]{\small Worksheet \##2}%
  \fancyfoot[L]{\small \SpeedExamLine}%
  \fancyfoot[C]{\small\thepage}%
  \fancyfoot[R]{\small \SpeedCredit}%
  \renewcommand{\footrulewidth}{0.4pt}%
}

% ── Title block (used at the top of every sheet AND answer file) ───────────────
% Usage: \SpeedTitleBlock{Daily Algebra Drill \#3}{Jane Doe}
% %% AGENTS: When generating a new sheet, ask the human contributor for the name they want credited in argument 2.
% For answer files, pass the "--- Answers \& Investigations" suffix in the arg.
\newcommand{\SpeedTitleBlock}[2]{%
  \begin{center}
    {\LARGE\bfseries #1}\\[4pt]
    {\normalsize\itshape Sheet by #2}\\[6pt]
    \hrule\vspace{4pt}
    \begin{tabular}{llcc}
      \toprule
      \textbf{Section} & \textbf{Topic} & \textbf{Questions} & \textbf{Target time}\\
      \midrule
      A & Rapid Recognition          & 10 & 2 min 30 sec \\
      B & Manipulation Drills        & 10 & 8 min        \\
      C & Substitution \& Structure  & 8  & 10 min       \\
      D & Challenge                  & 5  & 15 min       \\
      \bottomrule
    \end{tabular}
    \vspace{4pt}
    \hrule
  \end{center}
}

% ── Sheet metadata: topic + the day's new tools ────────────────────────────────
% Usage (immediately after \SpeedTitleBlock, sheets only):
%   \SpeedMeta{Expanding, factorising, and the identity toolkit}
%             {difference of squares, sum/product form, completing the square}
%
% Arg 1 = the sheet's topic in one line. The website lists this instead of
%         "Sheet 03", so write the thing a reader would actually search for.
% Arg 2 = comma-separated, at most five: the tools this sheet introduces that
%         earlier sheets in the pillar did not. These become the website's tags.
%         Leave out anything the reader already met — the tags are meant to say
%         what is new today, not to summarise the sheet.
%
% Both are parsed straight out of this file by tools/build_website.py, so the
% sheet stays the single source of truth and the site cannot drift from it.
%
% This renders NOTHING. It is a declaration for the build script, not a piece of
% the sheet's layout: adding it to a sheet must not change that sheet's PDF, so
% the 25 existing sheets can be annotated without a single one of them being
% reissued. If the toolkit should also appear on the page, that is a separate
% decision about the sheet design — make it deliberately, here, not by leaving
% this macro printing something nobody asked it to.
\newcommand{\SpeedMeta}[2]{}

% ── Answer / Method / Investigate-further boxes (answer files only) ────────────
\newmdenv[
  backgroundcolor=lightgray,
  linecolor=medgray,
  linewidth=0.5pt,
  leftmargin=0pt, rightmargin=0pt,
  innerleftmargin=8pt, innerrightmargin=8pt,
  innertopmargin=5pt, innerbottommargin=5pt,
  skipabove=4pt, skipbelow=4pt
]{ansbox}

\newmdenv[
  backgroundcolor=white,
  linecolor=darkgray,
  linewidth=0.8pt,
  leftline=true, rightline=false, topline=false, bottomline=false,
  leftmargin=0pt, rightmargin=0pt,
  innerleftmargin=8pt, innerrightmargin=6pt,
  innertopmargin=4pt, innerbottommargin=4pt,
  skipabove=4pt, skipbelow=6pt
]{invbox}

\newcommand{\ans}[1]{%
  \begin{ansbox}
    \textbf{Answer:} #1
  \end{ansbox}}

\newcommand{\method}[1]{%
  {\small\textbf{Method:} #1}\par}

\newcommand{\inv}[1]{%
  \begin{invbox}
    \textbf{\small Investigate further:} {\small #1}
  \end{invbox}}

% ── Misc helpers ────────────────────────────────────────────────────────────────
\newcommand{\qn}[1]{\textbf{#1.}\;}

% Mistake-linking: attach a Top 5 Patterns / Common Traps bullet to the question
% label(s) in this sheet where it actually shows up, e.g. \seealso{B6, D2}
\newcommand{\seealso}[1]{\hfill\textit{\footnotesize(see #1)}}

% ── Graph options macro ─────────────────────────────────────────────────────────
% Usage: \GraphOptions{tikz code for A}{tikz code for B}{tikz code for C}{tikz code for D}
\newcommand{\GraphOptions}[4]{%
  \begin{center}
    \begin{tabular}{cc}
      \textbf{A)} & \textbf{B)} \\
      #1 & #2 \\[10pt]
      \textbf{C)} & \textbf{D)} \\
      #3 & #4
    \end{tabular}
  \end{center}
}

% ── Closing motivational rule + cross-promo footer (sheets only) ────────────────
% Usage: \SpeedClosing{"Quote text."}
\newcommand{\SpeedClosing}[1]{%
  \vspace{20pt}
  \begin{center}
  \rule{0.6\textwidth}{0.4pt}\\[4pt]
  {\small\itshape #1}\\[4pt]
  \rule{0.6\textwidth}{0.4pt}
  \end{center}
  \begin{center}
  {\small Found a mistake or want to contribute a sheet?}\\
  {\small Visit the open-source project at \href{https://github.com/The-CerealDev/speed-maths}{github.com/The-CerealDev/speed-maths}}
  \end{center}
}

% Editing THIS file changes the look of every sheet/answer across every pillar.
% ══════════════════════════════════════════════════════════════════════════════

\usepackage[a4paper, top=25mm, bottom=25mm, left=22mm, right=22mm]{geometry}
\usepackage{amsmath, amssymb}
\usepackage{enumitem}
\usepackage{booktabs}
\usepackage{array}
\usepackage{parskip}
\usepackage{microtype}
\usepackage{fancyhdr}
\usepackage{titlesec}
\usepackage{mdframed}
\usepackage{xcolor}
\usepackage[hidelinks]{hyperref}
\usepackage{graphicx}
\usepackage{tikz}
\usepackage{pgfplots}
\pgfplotsset{compat=1.18}

% ── Colours ───────────────────────────────────────────────────────────────────
\definecolor{darkgray}{gray}{0.15}
\definecolor{medgray}{gray}{0.45}
\definecolor{lightgray}{gray}{0.93}
\definecolor{boxgray}{gray}{0.88}

% ── Section formatting ──────────────────────────────────────────────────────────
\titleformat{\section}[block]
  {\normalfont\large\bfseries}{}
  {0pt}{}[\vspace{2pt}\hrule\vspace{6pt}]
\titlespacing*{\section}{0pt}{14pt}{4pt}

% ── List formatting ─────────────────────────────────────────────────────────────
\setlist[enumerate,1]{label=\textbf{\arabic*.}, leftmargin=2em, itemsep=10pt}

% ── Branding constants (edit here, not per-file) ────────────────────────────────
\newcommand{\SpeedExamLine}{\textbf{Speed Maths:} Competition \& Exam Prep}
\newcommand{\SpeedCredit}{Series by \href{https://www.linkedin.com/in/cerealdev/}{David Akinyele-Aje}}

% ── Header / footer ──────────────────────────────────────────────────────────────
% Usage (once, after % main (standalone preamble for editing)
% vim: nowrap: spell:
% ══════════════════════════════════════════════════════════════════════════════
% Speed Maths --- shared preamble
% Included by every sheet and answer file via:  \input{../../shared/preamble}
% Editing THIS file changes the look of every sheet/answer across every pillar.
% ══════════════════════════════════════════════════════════════════════════════

\usepackage[a4paper, top=25mm, bottom=25mm, left=22mm, right=22mm]{geometry}
\usepackage{amsmath, amssymb}
\usepackage{enumitem}
\usepackage{booktabs}
\usepackage{array}
\usepackage{parskip}
\usepackage{microtype}
\usepackage{fancyhdr}
\usepackage{titlesec}
\usepackage{mdframed}
\usepackage{xcolor}
\usepackage[hidelinks]{hyperref}
\usepackage{graphicx}
\usepackage{tikz}
\usepackage{pgfplots}
\pgfplotsset{compat=1.18}

% ── Colours ───────────────────────────────────────────────────────────────────
\definecolor{darkgray}{gray}{0.15}
\definecolor{medgray}{gray}{0.45}
\definecolor{lightgray}{gray}{0.93}
\definecolor{boxgray}{gray}{0.88}

% ── Section formatting ──────────────────────────────────────────────────────────
\titleformat{\section}[block]
  {\normalfont\large\bfseries}{}
  {0pt}{}[\vspace{2pt}\hrule\vspace{6pt}]
\titlespacing*{\section}{0pt}{14pt}{4pt}

% ── List formatting ─────────────────────────────────────────────────────────────
\setlist[enumerate,1]{label=\textbf{\arabic*.}, leftmargin=2em, itemsep=10pt}

% ── Branding constants (edit here, not per-file) ────────────────────────────────
\newcommand{\SpeedExamLine}{\textbf{Speed Maths:} Competition \& Exam Prep}
\newcommand{\SpeedCredit}{Series by \href{https://www.linkedin.com/in/cerealdev/}{David Akinyele-Aje}}

% ── Header / footer ──────────────────────────────────────────────────────────────
% Usage (once, after \input{../../shared/preamble} and before \begin{document}):
%   \SpeedHeader{Algebra}{3}
% First arg  = pillar name shown as "Speed <Pillar> --- Daily Drill"
% Second arg = worksheet number
\pagestyle{fancy}
\newcommand{\SpeedHeader}[2]{%
  \fancyhf{}%
  \renewcommand{\headrulewidth}{0.4pt}%
  \setlength{\headheight}{14pt}%
  \fancyhead[L]{\small\textbf{Speed #1 --- Daily Drill}}%
  \fancyhead[R]{\small Worksheet \##2}%
  \fancyfoot[L]{\small \SpeedExamLine}%
  \fancyfoot[C]{\small\thepage}%
  \fancyfoot[R]{\small \SpeedCredit}%
  \renewcommand{\footrulewidth}{0.4pt}%
}

% ── Title block (used at the top of every sheet AND answer file) ───────────────
% Usage: \SpeedTitleBlock{Daily Algebra Drill \#3}{Jane Doe}
% %% AGENTS: When generating a new sheet, ask the human contributor for the name they want credited in argument 2.
% For answer files, pass the "--- Answers \& Investigations" suffix in the arg.
\newcommand{\SpeedTitleBlock}[2]{%
  \begin{center}
    {\LARGE\bfseries #1}\\[4pt]
    {\normalsize\itshape Sheet by #2}\\[6pt]
    \hrule\vspace{4pt}
    \begin{tabular}{llcc}
      \toprule
      \textbf{Section} & \textbf{Topic} & \textbf{Questions} & \textbf{Target time}\\
      \midrule
      A & Rapid Recognition          & 10 & 2 min 30 sec \\
      B & Manipulation Drills        & 10 & 8 min        \\
      C & Substitution \& Structure  & 8  & 10 min       \\
      D & Challenge                  & 5  & 15 min       \\
      \bottomrule
    \end{tabular}
    \vspace{4pt}
    \hrule
  \end{center}
}

% ── Sheet metadata: topic + the day's new tools ────────────────────────────────
% Usage (immediately after \SpeedTitleBlock, sheets only):
%   \SpeedMeta{Expanding, factorising, and the identity toolkit}
%             {difference of squares, sum/product form, completing the square}
%
% Arg 1 = the sheet's topic in one line. The website lists this instead of
%         "Sheet 03", so write the thing a reader would actually search for.
% Arg 2 = comma-separated, at most five: the tools this sheet introduces that
%         earlier sheets in the pillar did not. These become the website's tags.
%         Leave out anything the reader already met — the tags are meant to say
%         what is new today, not to summarise the sheet.
%
% Both are parsed straight out of this file by tools/build_website.py, so the
% sheet stays the single source of truth and the site cannot drift from it.
%
% This renders NOTHING. It is a declaration for the build script, not a piece of
% the sheet's layout: adding it to a sheet must not change that sheet's PDF, so
% the 25 existing sheets can be annotated without a single one of them being
% reissued. If the toolkit should also appear on the page, that is a separate
% decision about the sheet design — make it deliberately, here, not by leaving
% this macro printing something nobody asked it to.
\newcommand{\SpeedMeta}[2]{}

% ── Answer / Method / Investigate-further boxes (answer files only) ────────────
\newmdenv[
  backgroundcolor=lightgray,
  linecolor=medgray,
  linewidth=0.5pt,
  leftmargin=0pt, rightmargin=0pt,
  innerleftmargin=8pt, innerrightmargin=8pt,
  innertopmargin=5pt, innerbottommargin=5pt,
  skipabove=4pt, skipbelow=4pt
]{ansbox}

\newmdenv[
  backgroundcolor=white,
  linecolor=darkgray,
  linewidth=0.8pt,
  leftline=true, rightline=false, topline=false, bottomline=false,
  leftmargin=0pt, rightmargin=0pt,
  innerleftmargin=8pt, innerrightmargin=6pt,
  innertopmargin=4pt, innerbottommargin=4pt,
  skipabove=4pt, skipbelow=6pt
]{invbox}

\newcommand{\ans}[1]{%
  \begin{ansbox}
    \textbf{Answer:} #1
  \end{ansbox}}

\newcommand{\method}[1]{%
  {\small\textbf{Method:} #1}\par}

\newcommand{\inv}[1]{%
  \begin{invbox}
    \textbf{\small Investigate further:} {\small #1}
  \end{invbox}}

% ── Misc helpers ────────────────────────────────────────────────────────────────
\newcommand{\qn}[1]{\textbf{#1.}\;}

% Mistake-linking: attach a Top 5 Patterns / Common Traps bullet to the question
% label(s) in this sheet where it actually shows up, e.g. \seealso{B6, D2}
\newcommand{\seealso}[1]{\hfill\textit{\footnotesize(see #1)}}

% ── Graph options macro ─────────────────────────────────────────────────────────
% Usage: \GraphOptions{tikz code for A}{tikz code for B}{tikz code for C}{tikz code for D}
\newcommand{\GraphOptions}[4]{%
  \begin{center}
    \begin{tabular}{cc}
      \textbf{A)} & \textbf{B)} \\
      #1 & #2 \\[10pt]
      \textbf{C)} & \textbf{D)} \\
      #3 & #4
    \end{tabular}
  \end{center}
}

% ── Closing motivational rule + cross-promo footer (sheets only) ────────────────
% Usage: \SpeedClosing{"Quote text."}
\newcommand{\SpeedClosing}[1]{%
  \vspace{20pt}
  \begin{center}
  \rule{0.6\textwidth}{0.4pt}\\[4pt]
  {\small\itshape #1}\\[4pt]
  \rule{0.6\textwidth}{0.4pt}
  \end{center}
  \begin{center}
  {\small Found a mistake or want to contribute a sheet?}\\
  {\small Visit the open-source project at \href{https://github.com/The-CerealDev/speed-maths}{github.com/The-CerealDev/speed-maths}}
  \end{center}
}
 and before \begin{document}):
%   \SpeedHeader{Algebra}{3}
% First arg  = pillar name shown as "Speed <Pillar> --- Daily Drill"
% Second arg = worksheet number
\pagestyle{fancy}
\newcommand{\SpeedHeader}[2]{%
  \fancyhf{}%
  \renewcommand{\headrulewidth}{0.4pt}%
  \setlength{\headheight}{14pt}%
  \fancyhead[L]{\small\textbf{Speed #1 --- Daily Drill}}%
  \fancyhead[R]{\small Worksheet \##2}%
  \fancyfoot[L]{\small \SpeedExamLine}%
  \fancyfoot[C]{\small\thepage}%
  \fancyfoot[R]{\small \SpeedCredit}%
  \renewcommand{\footrulewidth}{0.4pt}%
}

% ── Title block (used at the top of every sheet AND answer file) ───────────────
% Usage: \SpeedTitleBlock{Daily Algebra Drill \#3}{Jane Doe}
% %% AGENTS: When generating a new sheet, ask the human contributor for the name they want credited in argument 2.
% For answer files, pass the "--- Answers \& Investigations" suffix in the arg.
\newcommand{\SpeedTitleBlock}[2]{%
  \begin{center}
    {\LARGE\bfseries #1}\\[4pt]
    {\normalsize\itshape Sheet by #2}\\[6pt]
    \hrule\vspace{4pt}
    \begin{tabular}{llcc}
      \toprule
      \textbf{Section} & \textbf{Topic} & \textbf{Questions} & \textbf{Target time}\\
      \midrule
      A & Rapid Recognition          & 10 & 2 min 30 sec \\
      B & Manipulation Drills        & 10 & 8 min        \\
      C & Substitution \& Structure  & 8  & 10 min       \\
      D & Challenge                  & 5  & 15 min       \\
      \bottomrule
    \end{tabular}
    \vspace{4pt}
    \hrule
  \end{center}
}

% ── Sheet metadata: topic + the day's new tools ────────────────────────────────
% Usage (immediately after \SpeedTitleBlock, sheets only):
%   \SpeedMeta{Expanding, factorising, and the identity toolkit}
%             {difference of squares, sum/product form, completing the square}
%
% Arg 1 = the sheet's topic in one line. The website lists this instead of
%         "Sheet 03", so write the thing a reader would actually search for.
% Arg 2 = comma-separated, at most five: the tools this sheet introduces that
%         earlier sheets in the pillar did not. These become the website's tags.
%         Leave out anything the reader already met — the tags are meant to say
%         what is new today, not to summarise the sheet.
%
% Both are parsed straight out of this file by tools/build_website.py, so the
% sheet stays the single source of truth and the site cannot drift from it.
%
% This renders NOTHING. It is a declaration for the build script, not a piece of
% the sheet's layout: adding it to a sheet must not change that sheet's PDF, so
% the 25 existing sheets can be annotated without a single one of them being
% reissued. If the toolkit should also appear on the page, that is a separate
% decision about the sheet design — make it deliberately, here, not by leaving
% this macro printing something nobody asked it to.
\newcommand{\SpeedMeta}[2]{}

% ── Answer / Method / Investigate-further boxes (answer files only) ────────────
\newmdenv[
  backgroundcolor=lightgray,
  linecolor=medgray,
  linewidth=0.5pt,
  leftmargin=0pt, rightmargin=0pt,
  innerleftmargin=8pt, innerrightmargin=8pt,
  innertopmargin=5pt, innerbottommargin=5pt,
  skipabove=4pt, skipbelow=4pt
]{ansbox}

\newmdenv[
  backgroundcolor=white,
  linecolor=darkgray,
  linewidth=0.8pt,
  leftline=true, rightline=false, topline=false, bottomline=false,
  leftmargin=0pt, rightmargin=0pt,
  innerleftmargin=8pt, innerrightmargin=6pt,
  innertopmargin=4pt, innerbottommargin=4pt,
  skipabove=4pt, skipbelow=6pt
]{invbox}

\newcommand{\ans}[1]{%
  \begin{ansbox}
    \textbf{Answer:} #1
  \end{ansbox}}

\newcommand{\method}[1]{%
  {\small\textbf{Method:} #1}\par}

\newcommand{\inv}[1]{%
  \begin{invbox}
    \textbf{\small Investigate further:} {\small #1}
  \end{invbox}}

% ── Misc helpers ────────────────────────────────────────────────────────────────
\newcommand{\qn}[1]{\textbf{#1.}\;}

% Mistake-linking: attach a Top 5 Patterns / Common Traps bullet to the question
% label(s) in this sheet where it actually shows up, e.g. \seealso{B6, D2}
\newcommand{\seealso}[1]{\hfill\textit{\footnotesize(see #1)}}

% ── Graph options macro ─────────────────────────────────────────────────────────
% Usage: \GraphOptions{tikz code for A}{tikz code for B}{tikz code for C}{tikz code for D}
\newcommand{\GraphOptions}[4]{%
  \begin{center}
    \begin{tabular}{cc}
      \textbf{A)} & \textbf{B)} \\
      #1 & #2 \\[10pt]
      \textbf{C)} & \textbf{D)} \\
      #3 & #4
    \end{tabular}
  \end{center}
}

% ── Closing motivational rule + cross-promo footer (sheets only) ────────────────
% Usage: \SpeedClosing{"Quote text."}
\newcommand{\SpeedClosing}[1]{%
  \vspace{20pt}
  \begin{center}
  \rule{0.6\textwidth}{0.4pt}\\[4pt]
  {\small\itshape #1}\\[4pt]
  \rule{0.6\textwidth}{0.4pt}
  \end{center}
  \begin{center}
  {\small Found a mistake or want to contribute a sheet?}\\
  {\small Visit the open-source project at \href{https://github.com/The-CerealDev/speed-maths}{github.com/The-CerealDev/speed-maths}}
  \end{center}
}
 and before \begin{document}):
%   \SpeedHeader{Algebra}{3}
% First arg  = pillar name shown as "Speed <Pillar> --- Daily Drill"
% Second arg = worksheet number
\pagestyle{fancy}
\newcommand{\SpeedHeader}[2]{%
  \fancyhf{}%
  \renewcommand{\headrulewidth}{0.4pt}%
  \setlength{\headheight}{14pt}%
  \fancyhead[L]{\small\textbf{Speed #1 --- Daily Drill}}%
  \fancyhead[R]{\small Worksheet \##2}%
  \fancyfoot[L]{\small \SpeedExamLine}%
  \fancyfoot[C]{\small\thepage}%
  \fancyfoot[R]{\small \SpeedCredit}%
  \renewcommand{\footrulewidth}{0.4pt}%
}

% ── Title block (used at the top of every sheet AND answer file) ───────────────
% Usage: \SpeedTitleBlock{Daily Algebra Drill \#3}{Jane Doe}
% %% AGENTS: When generating a new sheet, ask the human contributor for the name they want credited in argument 2.
% For answer files, pass the "--- Answers \& Investigations" suffix in the arg.
\newcommand{\SpeedTitleBlock}[2]{%
  \begin{center}
    {\LARGE\bfseries #1}\\[4pt]
    {\normalsize\itshape Sheet by #2}\\[6pt]
    \hrule\vspace{4pt}
    \begin{tabular}{llcc}
      \toprule
      \textbf{Section} & \textbf{Topic} & \textbf{Questions} & \textbf{Target time}\\
      \midrule
      A & Rapid Recognition          & 10 & 2 min 30 sec \\
      B & Manipulation Drills        & 10 & 8 min        \\
      C & Substitution \& Structure  & 8  & 10 min       \\
      D & Challenge                  & 5  & 15 min       \\
      \bottomrule
    \end{tabular}
    \vspace{4pt}
    \hrule
  \end{center}
}

% ── Sheet metadata: topic + the day's new tools ────────────────────────────────
% Usage (immediately after \SpeedTitleBlock, sheets only):
%   \SpeedMeta{Expanding, factorising, and the identity toolkit}
%             {difference of squares, sum/product form, completing the square}
%
% Arg 1 = the sheet's topic in one line. The website lists this instead of
%         "Sheet 03", so write the thing a reader would actually search for.
% Arg 2 = comma-separated, at most five: the tools this sheet introduces that
%         earlier sheets in the pillar did not. These become the website's tags.
%         Leave out anything the reader already met — the tags are meant to say
%         what is new today, not to summarise the sheet.
%
% Both are parsed straight out of this file by tools/build_website.py, so the
% sheet stays the single source of truth and the site cannot drift from it.
%
% This renders NOTHING. It is a declaration for the build script, not a piece of
% the sheet's layout: adding it to a sheet must not change that sheet's PDF, so
% the 25 existing sheets can be annotated without a single one of them being
% reissued. If the toolkit should also appear on the page, that is a separate
% decision about the sheet design — make it deliberately, here, not by leaving
% this macro printing something nobody asked it to.
\newcommand{\SpeedMeta}[2]{}

% ── Answer / Method / Investigate-further boxes (answer files only) ────────────
\newmdenv[
  backgroundcolor=lightgray,
  linecolor=medgray,
  linewidth=0.5pt,
  leftmargin=0pt, rightmargin=0pt,
  innerleftmargin=8pt, innerrightmargin=8pt,
  innertopmargin=5pt, innerbottommargin=5pt,
  skipabove=4pt, skipbelow=4pt
]{ansbox}

\newmdenv[
  backgroundcolor=white,
  linecolor=darkgray,
  linewidth=0.8pt,
  leftline=true, rightline=false, topline=false, bottomline=false,
  leftmargin=0pt, rightmargin=0pt,
  innerleftmargin=8pt, innerrightmargin=6pt,
  innertopmargin=4pt, innerbottommargin=4pt,
  skipabove=4pt, skipbelow=6pt
]{invbox}

\newcommand{\ans}[1]{%
  \begin{ansbox}
    \textbf{Answer:} #1
  \end{ansbox}}

\newcommand{\method}[1]{%
  {\small\textbf{Method:} #1}\par}

\newcommand{\inv}[1]{%
  \begin{invbox}
    \textbf{\small Investigate further:} {\small #1}
  \end{invbox}}

% ── Misc helpers ────────────────────────────────────────────────────────────────
\newcommand{\qn}[1]{\textbf{#1.}\;}

% Mistake-linking: attach a Top 5 Patterns / Common Traps bullet to the question
% label(s) in this sheet where it actually shows up, e.g. \seealso{B6, D2}
\newcommand{\seealso}[1]{\hfill\textit{\footnotesize(see #1)}}

% ── Graph options macro ─────────────────────────────────────────────────────────
% Usage: \GraphOptions{tikz code for A}{tikz code for B}{tikz code for C}{tikz code for D}
\newcommand{\GraphOptions}[4]{%
  \begin{center}
    \begin{tabular}{cc}
      \textbf{A)} & \textbf{B)} \\
      #1 & #2 \\[10pt]
      \textbf{C)} & \textbf{D)} \\
      #3 & #4
    \end{tabular}
  \end{center}
}

% ── Closing motivational rule + cross-promo footer (sheets only) ────────────────
% Usage: \SpeedClosing{"Quote text."}
\newcommand{\SpeedClosing}[1]{%
  \vspace{20pt}
  \begin{center}
  \rule{0.6\textwidth}{0.4pt}\\[4pt]
  {\small\itshape #1}\\[4pt]
  \rule{0.6\textwidth}{0.4pt}
  \end{center}
  \begin{center}
  {\small Found a mistake or want to contribute a sheet?}\\
  {\small Visit the open-source project at \href{https://github.com/The-CerealDev/speed-maths}{github.com/The-CerealDev/speed-maths}}
  \end{center}
}

\SpeedHeader{Geometry}{1}

\begin{document}

\SpeedTitleBlock{Daily Geometry Drill \#1 --- Answers \& Investigations}{Henry Koduthore}
\noindent\textit{New toolkit today: Gradients \& Perpendicularity $\cdot$ Distance \& Midpoints $\cdot$ Perpendicular Distance $\cdot$ Shoelace Areas $\cdot$ Collinearity.}

%═══════════════════════════════════════════════════════════════════════════════
\section*{Section A \quad Rapid Recognition}
%═══════════════════════════════════════════════════════════════════════════════
\begin{enumerate}

%── A1 ──────────────────────────────────────────────────────────────────────────
\item Find the gradient of the line $\;6x - 3y + 15 = 0$.

\ans{$2$}
\method{Rearrange to gradient--intercept form: $3y=6x+15\Rightarrow y=2x+5$, so $m=2$.}
\inv{Implicit form $ax+by+c=0$ has gradient $-a/b$; here $-\frac{6}{-3}=2$. \textbf{If the line is instead written $3y-6x-15=0$ (the same line), is the gradient still $2$?} It must be --- check the $-a/b$ rule still agrees.}

%── A2 ──────────────────────────────────────────────────────────────────────────
\item Find the gradient of the line through $(1,2)$ and $(5,14)$.

\ans{$3$}
\method{$m=\frac{14-2}{5-1}=\frac{12}{4}=3$.}
\inv{Check with $(5-1,14-2)$ swapped: $\frac{2-14}{1-5}=\frac{-12}{-4}=3$ \checkmark --- the sign always cancels. \textbf{If the line through $(1,2)$ and $(5,14)$ is moved up by $3$ (every $y$ increased by $3$), does its gradient change?} No --- a vertical translation keeps the ratio of rises fixed.}

%── A3 ──────────────────────────────────────────────────────────────────────────
\item A line of gradient $4$ is perpendicular to a line of gradient $m$. Write down $m$.

\ans{$-\frac{1}{4}$}
\method{Perpendicular gradients are negative reciprocals: $m=-\frac{1}{4}$, since $4\cdot(-\frac14)=-1$.}
\inv{The product test $m_1m_2=-1$ is the fastest check. \textbf{Is the line $y=4x$ parallel, perpendicular, or neither, to $y=-\frac14x+12$?} The constant term is a distraction --- gradients alone decide.}

%── A4 ──────────────────────────────────────────────────────────────────────────
\item Find the distance between $(2,1)$ and $(5,5)$.

\ans{$5$}
\method{$d=\sqrt{(5-2)^2+(5-1)^2}=\sqrt{9+16}=\sqrt{25}=5$.}
\inv{$3$-$4$-$5$ right triangle in disguise. \textbf{The points $(-1,0)$ and $(3,-3)$ are the same distance apart --- can you see the Pythagorean triple before computing, by drawing the displacement vector?} The displacement $(4,-3)$ is a reordered $3$-$4$-$5$.}

%── A5 ──────────────────────────────────────────────────────────────────────────
\item Find the midpoint of the segment joining $(4,8)$ and $(10,2)$.

\ans{"$(7,5)$"}
\method{Midpoint averages each coordinate: $\left(\frac{4+10}{2},\frac{8+2}{2}\right)=(7,5)$.}
\inv{Check both halves are equal distances: $(7-4,8-5)=(3,3)$ and $(10-7,5-2)=(3,3)$ \checkmark. \textbf{The midpoint of two points is the {\em average} of the two displacement vectors. What is the midpoint of $(4,8)$ and $(-6,-2)$?} Negative coordinates behave the same way.}

%── A6 ──────────────────────────────────────────────────────────────────────────
\item Find the perpendicular distance from the origin to the line $\;3x + 4y - 10 = 0$.

\ans{$2$}
\method{Use $\frac{|ax_0+by_0+c|}{\sqrt{a^2+b^2}}$ at $(0,0)$: $\frac{|-10|}{\sqrt{9+16}}=\frac{10}{5}=2$.}
\inv{The denominator $5$ is the length of the normal vector $(3,4)$. \textbf{The point of the line closest to the origin satisfies $(x,y)=t(3,4)$ for some $t$ --- find it and confirm its distance is indeed $2$.} Every point of the line's coordinate vector is a multiple of its normal.}

%── A7 ──────────────────────────────────────────────────────────────────────────
\item A line has $x$-intercept $4$ and $y$-intercept $6$. Find its gradient.

\ans{$-\frac{3}{2}$}
\method{The intercepts are the points $(4,0)$ and $(0,6)$: $m=\frac{6-0}{0-4}=-\frac{3}{2}$.}
\inv{Reading right-to-left: this is the negative of $\frac{\text{$y$-intercept}}{\text{$x$-intercept}}$. \textbf{Write the line in intercept form $\frac{x}{4}+\frac{y}{6}=1$ and rearrange to confirm $m=-\frac{y\text{-intercept}}{x\text{-intercept}}=-\frac{6}{4}=-\frac32$ --- the intercept form makes the gradient transparent.}}

%── A8 ──────────────────────────────────────────────────────────────────────────
\item Find the area of the triangle with vertices $(0,0)$, $(4,0)$ and $(0,3)$.

\ans{$6$}
\method{Right-angled at the origin: area $=\frac12\cdot4\cdot3=6$.}
\inv{Equivalently the shoelace/midpoint-vector formula $\frac12|(4)(3)-(0)(0)|=6$. \textbf{Rotate this triangle by $90^\circ$ about the origin: does the area change, and do you get the same shoelace number?} Reflections and rotations preserve area.}

%── A9 ──────────────────────────────────────────────────────────────────────────
\item The points $(1,2)$, $(2,4)$ and $(k,10)$ are collinear. Find $k$.

\ans{$5$}
\method{Rise per unit run is $2$ (from $(1,2)$ to $(2,4)$); continuing to $y=10$ climbs $6$ more, so $k=2+\frac{6}{2}=5$.}
\inv{Formal check: $\frac{4-2}{2-1}=\frac{10-4}{k-2}\Rightarrow 2=\frac{6}{k-2}\Rightarrow k=5$. \textbf{What happens to $k$ if the points are required to be collinear but the order is $(2,4)$, $(1,2)$, $(k,10)$?} Collinearity is order-independent --- all three gradients must agree.}

%── A10 ─────────────────────────────────────────────────────────────────────────
\item Using the shoelace formula, find the area of the triangle with vertices $(0,0)$, $(3,4)$ and $(6,0)$.

\ans{$12$}
\method{Shoelace: $\frac12|0\cdot4+3\cdot0+6\cdot0-(0\cdot3+4\cdot6+0\cdot0)|=\frac12|24|=12$.}
\inv{Bonus route: base $=6$, height $=4$, area $=\frac12\cdot6\cdot4=12$ \checkmark. \textbf{For a triangle with one vertex at the origin, shoelace collapses to $\frac12|x_1y_2-x_2y_1|$ with the other two points --- prove that this equals $\frac12|\det|$ of the two vectors.} That determinant is the signed area.}

\end{enumerate}

%═══════════════════════════════════════════════════════════════════════════════
\section*{Section B \quad Manipulation Drills}
%═══════════════════════════════════════════════════════════════════════════════
\begin{enumerate}

%── B1 ──────────────────────────────────────────────────────────────────────────
\item The line $\;2x - 3y + 6 = 0\;$ is perpendicular to which of the following?
\begin{itemize}
\item[A)] $y = -\dfrac{3}{2}x + 1$
\item[B)] $y = \dfrac{3}{2}x - 2$
\item[C)] $y = \dfrac{2}{3}x + 5$
\item[D)] $y = -\dfrac{2}{3}x + 3$
\end{itemize}

\ans{A) $y = -\frac{3}{2}x + 1$}
\method{Given line has gradient $\frac23$ (rule $-a/b$), so perpendicular gradient is $-\frac32$. Only option A has gradient $-\frac32$.}
\inv{Options B and C have the same gradient, and D the new ``same slope'' decoy --- read \emph{gradients only} through the $-\frac32$ filter. \textbf{Which option is parallel to the given line, and which is merely reflection-symmetric in slope?} Perpendicularity and parallelism are decided by one number, the slope.}

%── B2 ──────────────────────────────────────────────────────────────────────────
\item What is the perpendicular distance from $(2,3)$ to the line $\;5x + 12y - 7 = 0$?
\begin{itemize}
\item[A)] $1$
\item[B)] $2$
\item[C)] $3$
\item[D)] $4$
\end{itemize}

\ans{C) $3$}
\method{$\frac{|5(2)+12(3)-7|}{\sqrt{25+144}}=\frac{|10+36-7|}{13}=\frac{39}{13}=3$.}
\inv{The numerator is exactly how far the point sits from the line measured along a normal whose length is $13$. \textbf{Two points are at perpendicular distance $3$ from $5x+12y-7=0$ --- what do they have in common?} They lie on $5x+12y-7=\pm39$, the two parallel satellite lines.}

%── B3 ──────────────────────────────────────────────────────────────────────────
\item Which point lies on the perpendicular bisector of the segment joining $(1,2)$ and $(5,6)$?
\begin{itemize}
\item[A)] $(2,5)$
\item[B)] $(3,3)$
\item[C)] $(4,4)$
\item[D)] $(5,3)$
\end{itemize}

\ans{A) $(2,5)$}
\method{Midpoint $(3,4)$; segment gradient $1$, so bisector gradient $-1$: $y-4=-(x-3)$, i.e.\ $x+y=7$. Check options: $2+5=7$ \checkmark.}
\inv{Only A satisfies $x+y=7$. \textbf{A point is on the perpendicular bisector iff it is {\em equidistant} from the two endpoints --- verify $(2,5)$ is distance $\sqrt{16}$ from both $(1,2)$ and $(5,6)$.} That's the defining property, worth internalising.}

%── B4 ──────────────────────────────────────────────────────────────────────────
\item Find the area of the quadrilateral with vertices $(0,0)$, $(4,0)$, $(5,3)$, $(1,3)$, in that order.
\begin{itemize}
\item[A)] $10$
\item[B)] $12$
\item[C)] $14$
\item[D)] $16$
\end{itemize}

\ans{B) $12$}
\method{Shoelace in order: $\frac12|0\cdot0+4\cdot3+5\cdot3+1\cdot0-(0\cdot4+0\cdot5+3\cdot1+3\cdot0)|=\frac12|0+12+15+0-(0+0+3+0)|=\frac12|24|=12$.}
\inv{The quadrilateral is a parallelogram (bases $\parallel$ via position vectors $(4,0)$ and $(1,3)$), so area $=|\det|=|4(3)-0(1)|=12$ \checkmark. \textbf{Reorder the vertices as $(0,0),(5,3),(4,0),(1,3)$: the same region, but the shoelace number warns you when points are listed in the wrong order.} Keep vertices in cyclic order.}

%── B5 ──────────────────────────────────────────────────────────────────────────
\item The line through $(1,2)$ and $(4,5)$ crosses the $y$-axis at which value?
\begin{itemize}
\item[A)] $0$
\item[B)] $1$
\item[C)] $2$
\item[D)] $3$
\end{itemize}

\ans{B) $1$}
\method{Gradient $\frac{5-2}{4-1}=1$; $y=x+c$ through $(1,2)$ gives $c=1$.}
\inv{Extension: at $x=0$, $y=1$. \textbf{Through the point-slope form, $y-2=1(x-1)$. If the line were instead the one through $(1,2)$ and $(4,8)$, the $y$-intercept changes --- resolve $c$ again.} Always plug the known point in, never guess $c$ from the picture.}

%── B6 ──────────────────────────────────────────────────────────────────────────
\item Find the perpendicular distance from $(0,1)$ to the line $\;6x - 8y + 2 = 0$.

\ans{$\frac{3}{5}$}
\method{$\frac{|6(0)-8(1)+2|}{\sqrt{36+64}}=\frac{|-6|}{10}=\frac{3}{5}$.}
\inv{Normal length $\sqrt{36+64}=10$; numerator $6$. \textbf{The line $6x-8y+2=0$ sits distance $\frac15$ from the origin (check $\frac{|2|}{10}$). Since the $y$-axis point $(0,1)$ is $1$ unit up from the origin, the net distance $|1-\frac15|=\frac45$ differs from $\frac35$ --- why?} Because the normal $(6,-8)$ is not vertical: only the projection along it counts.}

%── B7 ──────────────────────────────────────────────────────────────────────────
\item Which of the following is NOT parallel to $\;y = 3x + 1$?
\begin{itemize}
\item[A)] $y = 3x + 5$
\item[B)] $6x - 2y = 4$
\item[C)] $3x - y = 2$
\item[D)] $y = 2x + 1$
\end{itemize}

\ans{D) $y = 2x + 1$}
\method{Parallel $\Leftrightarrow$ same gradient $3$. A: $3$\checkmark; B: $6x-2y=4\Rightarrow y=3x-2$\checkmark; C: $3x-y=2\Rightarrow y=3x-2$\checkmark; D: $y=2x+1$ has gradient $2$.}
\inv{Options B and C are disguises (divide by the $y$-coefficient first). \textbf{Show two distinct parallel lines never intersect, while any two non-parallel lines always do.} Parallelism is exactly ``same slope, different intercept''.}

%── B8 ──────────────────────────────────────────────────────────────────────────
\item The line $y = x + 1$ meets the line through $(0,3)$ perpendicular to it at the point $P$. Find the coordinates of $P$.

\ans{"$(1,2)$"}
\method{Through $(0,3)$ with gradient $-1$: $y=-x+3$. Intersect: $x+1=-x+3\Rightarrow x=1$, then $y=2$.}
\inv{The intersection of two perpendicular lines is the right-angle vertex of a right triangle. Verify $P=(1,2)$ lies on both lines: $y=x+1$ gives $2=1+1$ \checkmark\ and $y=-x+3$ gives $2=-1+3$ \checkmark. \textbf{Through $(0,3)$ draw the perpendicular and reflect the point where it meets $y=x+1$ across that meeting point --- what do you get, and how does it connect to the midpoint of $(0,3)$ and $(2,1)$?} Perpendicular lines meet exactly once; the reflect-and-extend trick answers the follow-up.}

%── B9 ──────────────────────────────────────────────────────────────────────────
\item Which point is closest to the origin?
\begin{itemize}
\item[A)] $(1,3)$
\item[B)] $(2,2)$
\item[C)] $(4,0)$
\item[D)] $(0,5)$
\end{itemize}

\ans{B) $(2,2)$}
\method{Distances: $\sqrt{(1,3)}=\sqrt{10}$, $\sqrt{(2,2)}=\sqrt{8}$, `$(4,0)$'$=4$, `$(0,5)$'$=5$; smallest is $\sqrt{8}$.}
\inv{Comparing squared distances avoids surds: $10$, $8$, $16$, $25$. \textbf{Among points with $x,y$ nonnegative integers and $x+y=4$, which is closest to the origin?} Both squares and sum compare against $(2,2)$.}

%── B10 ─────────────────────────────────────────────────────────────────────────
\item Using the shoelace formula, find the area of the triangle with vertices $(1,1)$, $(4,2)$, $(2,6)$.

\ans{$7$}
\method{$\frac12|1\cdot2+4\cdot6+2\cdot1-(1\cdot4+2\cdot2+6\cdot1)|=\frac12|2+24+2-(4+4+6)|=\frac12|28-14|=7$.}
\inv{The $7$ here is exactly the determinant area of the two vectors $(3,1)$ and $(1,5)$: $\frac12|3(5)-1(1)|=\frac12(14)=7$ \checkmark. \textbf{In shoelace, matching -- and cancelling -- the repeated coordinate arithmetic is a good fastest-check habit.} Confirm you did $\frac12|(x_1y_2+x_2y_3+x_3y_1)-(y_1x_2+y_2x_3+y_3x_1)|$.}

\end{enumerate}

%═══════════════════════════════════════════════════════════════════════════════
\section*{Section C \quad Substitution \& Structure}
%═══════════════════════════════════════════════════════════════════════════════
\begin{enumerate}

%── C1 ──────────────────────────────────────────────────────────────────────────
\item The midpoint of $(2,-1)$ and $(6,7)$ is?
\begin{itemize}
\item[A)] $(4,3)$
\item[B)] $(3,4)$
\item[C)] $(5,4)$
\item[D)] $(4,2)$
\end{itemize}

\ans{A) $(4,3)$}
\method{$\left(\frac{2+6}{2},\frac{-1+7}{2}\right)=(4,3)$.}
\inv{A and B swap the coordinates of the answer; D averages the $x$'s but not the $y$'s. \textbf{Midpoint of $(2,-1)$ and $(6,7)$ is {(4,3)}; of $(2,-1)$ and $(6,7)$ reflected in the origin, $(-6,-7)$, is $(-2,-4)$.} Averaging negatives works the same.}

%── C2 ──────────────────────────────────────────────────────────────────────────
\item The perpendicular bisector of the segment joining $(0,1)$ and $(4,5)$ has equation?
\begin{itemize}
\item[A)] $x + y = 5$
\item[B)] $y = x + 1$
\item[C)] $y = -x + 3$
\item[D)] $y = x + 3$
\end{itemize}

\ans{A) $x + y = 5$}
\method{Midpoint $(2,3)$; segment gradient $\frac{5-1}{4-0}=1$, so perpendicular gradient $-1$: $y-3=-(x-2)\Rightarrow x+y=5$.}
\inv{B and D are the perpendicular-doubles (gradient $1$); C is a parallel shifted copy missing the midpoint. \textbf{The perpendicular bisector is the locus of equidistant points --- verify $(2,3)$ and, say, $(6,-1)$ both satisfy $x+y=5$ and check equal distances to both endpoints.}}

%── C3 ──────────────────────────────────────────────────────────────────────────
\item The distance between $(2,5)$ and $(7,17)$ is?
\begin{itemize}
\item[A)] $13$
\item[B)] $14$
\item[C)] $\sqrt{194}$
\item[D)] $12$
\end{itemize}

\ans{A) $13$}
\method{$\sqrt{(7-2)^2+(17-5)^2}=\sqrt{25+144}=\sqrt{169}=13$.}
\inv{5-12-13 triple again; option C is the ``forgot to take the square root'' trap ($25+144=169$, not $\sqrt{194}$). \textbf{If the second point were $(7,13)$ instead, the distance $\sqrt{25+64}=\sqrt{89}$ stays a surd --- prove it is irrational.} Rational squares never land on $89$.}

%── C4 ──────────────────────────────────────────────────────────────────────────
\item The area of the triangle with vertices $(0,0)$, $(5,0)$, $(2,6)$ is?
\begin{itemize}
\item[A)] $15$
\item[B)] $30$
\item[C)] $10$
\item[D)] $12$
\end{itemize}

\ans{A) $15$}
\method{Take $(5,0)$ as base, height $6$ from $(2,6)$: $\frac12\cdot5\cdot6=15$.}
\inv{Shoelace check: $\frac12|0+5(6)+0-(0+0+0)|=15$ \checkmark. \textbf{The altitude is measured to the {\em base line} $y=0$ --- the $6$ is the $y$-coordinate. If the apex were $(2,-6)$ the area is still $15$; is that a reflection or a genuine second triangle?} Area is absolute.}

%── C5 ──────────────────────────────────────────────────────────────────────────
\item Which line is perpendicular to $\;y = 2x - 3$?
\begin{itemize}
\item[A)] $y = -2x + 1$
\item[B)] $y = -\frac{1}{2}x + 4$
\item[C)] $y = \frac{1}{2}x$
\item[D)] $y = 2x + 5$
\end{itemize}

\ans{B) $y = -\frac{1}{2}x + 4$}
\method{Gradient $2$ flips and inverts to $-\frac12$.}
\inv{A is the ``negate only'' trap, C the ``invert only'' trap, D parallel. \textbf{Which lines through $(0,4)$ are perpendicular to $y=2x-3$, and how many are there?} Through any fixed point there is exactly one perpendicular line: $y=-\tfrac12x+4$ here.}

%── C6 ──────────────────────────────────────────────────────────────────────────
\item The points $(1,4)$, $(4,7)$, $(10,k)$ are collinear. Find $k$.
\begin{itemize}
\item[A)] $10$
\item[B)] $12$
\item[C)] $13$
\item[D)] $15$
\end{itemize}

\ans{C) $13$}
\method{Gradient between the first two is $\frac{7-4}{4-1}=1$; at $x=10$, six units further, $y=4+\frac{10-1}{1}\cdot1=13$.}
\inv{Equivalently $k=4+(10-1)=13$, and $\frac{k-7}{10-4}=1\Rightarrow k=13$ \checkmark. \textbf{If the middle point were removed, the ``two'' gradients that must agree are $(1,4)\to(10,k)$ and $(1,4)\to(4,7)$ --- verify all three pairwise gradients equal $1$.}}

%── C7 ──────────────────────────────────────────────────────────────────────────
\item The perpendicular distance from $(2,-1)$ to the line $\;3x - 4y + 10 = 0$ is?
\begin{itemize}
\item[A)] $2$
\item[B)] $3$
\item[C)] $4$
\item[D)] $5$
\end{itemize}

\ans{C) $4$}
\method{$\frac{|3(2)-4(-1)+10|}{\sqrt{9+16}}=\frac{|6+4+10|}{5}=\frac{20}{5}=4$.}
\inv{Big numerator: the point lies far out along the normal. \textbf{Recompute with the point $(2,1)$: $\frac{|6-4+10|}{5}=\frac{12}{5}=2.4$; the distance changes as the point slides, while the normal length $5$ stays fixed.} The numerator measures signed offset along the normal.}

%── C8 ──────────────────────────────────────────────────────────────────────────
\item The line through $(2,1)$ perpendicular to $\;y = \frac{1}{2}x + 3$ has equation?
\begin{itemize}
\item[A)] $y = -2x + 5$
\item[B)] $y = 2x - 3$
\item[C)] $y = -\frac{1}{2}x + 1$
\item[D)] $y = -2x + 3$
\end{itemize}

\ans{A) $y = -2x + 5$}
\method{Gradient $-\frac{1}{1/2}=-2$; through $(2,1)$: $y-1=-2(x-2)\Rightarrow y=-2x+5$.}
\inv{Options A and D share gradient $-2$ but only A passes through $(2,1)$ (check $1=-4+5$ \checkmark). \textbf{The line $y=-2x+5$ and $y=\frac12x+3$ meet at $x=0.8,y=3.4$ --- verify the product of gradients $(-2)(\frac12)=-1$ \checkmark.} Perpendicular lines always cross at right angles, wherever that is.}

\end{enumerate}

%═══════════════════════════════════════════════════════════════════════════════
\section*{Section D \quad Challenge}
%═══════════════════════════════════════════════════════════════════════════════
\begin{enumerate}

%── D1 ──────────────────────────────────────────────────────────────────────────
\item The perpendicular bisector of the segment joining $(3,-1)$ and $(7,5)$ meets the $x$-axis at $P$. What is the $x$-coordinate of $P$?
\begin{itemize}
\item[A)] $6$
\item[B)] $7$
\item[C)] $8$
\item[D)] $9$
\end{itemize}

\ans{C) $8$}
\method{Midpoint $(5,2)$; segment gradient $\frac{5-(-1)}{7-3}=\frac32$, so bisector gradient $-\frac23$: $y-2=-\frac23(x-5)$. At $y=0$: $-2=-\frac23(x-5)\Rightarrow x-5=3\Rightarrow x=8$.}
\inv{General trick: the perpendicular bisector of a segment is equidistant from both ends, so $P=(p,0)$ satisfies $\sqrt{(p-3)^2+1}=\sqrt{(p-7)^2+25}$. Squaring: $(p-3)^2+1=(p-7)^2+25\Rightarrow p=8$ \checkmark. \textbf{Both routes agree --- the equidistance route never needs a slope.}}

%── D2 ──────────────────────────────────────────────────────────────────────────
\item A triangle with vertices $(0,0)$, $(6,0)$ and $(3,h)$ has area $12$. Find $h$.
\begin{itemize}
\item[A)] $2$
\item[B)] $4$
\item[C)] $6$
\item[D)] $8$
\end{itemize}

\ans{B) $4$}
\method{Base $6$ on the $x$-axis; altitude is $|h|$: $\frac12\cdot6\cdot|h|=12\Rightarrow h=\pm4$.}
\inv{Options B and D double the right answer? No --- B is $4$, D is $8$; both signs are valid, the question intends the positive one. \textbf{The apex $(3,h)$ is directly above the midpoint of the base, making an isosceles triangle. What is its perimeter at $h=4$?} Sides are $x=3,y=4$ hypotenuses of length $5$: perimeter $=6+10=16$.}

%── D3 ──────────────────────────────────────────────────────────────────────────
\item The distance from the point $(4,3)$ to the line $y = x$ is?
\begin{itemize}
\item[A)] $\dfrac{1}{\sqrt{2}}$
\item[B)] $1$
\item[C)] $\sqrt{2}$
\item[D)] $\dfrac{7}{\sqrt{2}}$
\end{itemize}

\ans{A) $\frac{1}{\sqrt{2}}$}
\method{Line $x-y=0$: $\frac{|4-3|}{\sqrt{1^2+(-1)^2}}=\frac{1}{\sqrt{2}}$.}
\inv{For $y=x$ the distance formula halves to $|x_0-y_0|/\sqrt{2}$ --- the $1$ difference over $\sqrt2$. \textbf{The point $(4,3)$ lies $\frac{1}{\sqrt{2}}$ from $y=x$ and $\frac{7}{\sqrt{2}}$ from $y=-x$ (check $|4+3|/\sqrt2$). What is the distance between the two lines $y=x$ and $y=-x$ themselves?} It's $0$ --- they meet at the origin; only non-parallel pairs have a ``distance'' in this sense, and these two are perpendicular.}

%── D4 ──────────────────────────────────────────────────────────────────────────
\item The line $\;2x + 3y - 12 = 0\;$ cuts the coordinate axes to form a right-angled triangle. Its area is?
\begin{itemize}
\item[A)] $6$
\item[B)] $8$
\item[C)] $12$
\item[D)] $24$
\end{itemize}

\ans{C) $12$}
\method{Intercepts: $x=6$ at $y=0$, $y=4$ at $x=0$, so area $=\frac12\cdot6\cdot4=12$.}
\inv{The intercept form $\frac{x}{6}+\frac{y}{4}=1$ reads intercepts straight off. \textbf{If instead the line were $4x+3y-12=0$, the intercepts are $3$ and $4$ --- which axis-touching area, $6$ or $12$, does the swap produce?} Check $\frac12\cdot3\cdot4=6$; the $12$ in the defining equation is {\em not} the area.}

%── D5 ──────────────────────────────────────────────────────────────────────────
\item Using the shoelace formula, find the area of the quadrilateral with vertices $(0,2)$, $(3,0)$, $(6,4)$, $(1,6)$, in that order.

\ans{$20$}
\method{Shoelace in order: $\frac12|(0)(0)+(3)(4)+(6)(6)+(1)(2)-[(2)(3)+(0)(6)+(4)(1)+(6)(0)]|=\frac12|(0+12+36+2)-(6+0+4+0)|=\frac12|50-10|=20$.}
\inv{Cross-check by triangulating along $(0,2)\!\to\!(6,4)$: the two triangles have areas $9$ and $11$, summing to $20$ \checkmark. \textbf{Now list the same four points in the order $(0,2),(3,0),(1,6),(6,4)$. Compute the shoelace value --- it differs, because that polygon crosses itself.} Reconcile the two, and note that area formulas assume a simple (non-crossing) polygon.}

\end{enumerate}

%═══════════════════════════════════════════════════════════════════════════════
\section*{Top 5 Patterns --- Worksheet \#1}
\begin{enumerate}[leftmargin=2em,itemsep=4pt]
  \item \textbf{Implicit line $\rightarrow$ gradient: rearrange to $y=mx+c$ (or read $-a/b$).} The constant $c$ never touches $m$ --- strip it off last. \seealso{A1, A7, B5}
  \item \textbf{Perpendicular gradients are negative reciprocals: $m_1m_2=-1$.} Slack jaw never flips and negates; keep both operations. \seealso{A3, B1, B8, C5, C8}
  \item \textbf{Perpendicular bisector $=$ midpoint $+$ negative-reciprocal slope.} Always hit the midpoint first; the bisector is the line through it with the flipped-inverted gradient. \seealso{B3, C2, D1}
  \item \textbf{Perpendicular distance formula: $\dfrac{|ax_0+by_0+c|}{\sqrt{a^2+b^2}}$.} Numerator from substitution, denominator from the normal vector's length. \seealso{A6, B2, B6, C7, D3}
  \item \textbf{Shoelace (signed) area: $\frac12\left|\sum x_iy_{i+1}-\sum y_ix_{i+1}\right|$.} Vertices in cyclic order; the $\frac12$ and the absolute value are non-negotiable. \seealso{A10, B4, B10, D5}
\end{enumerate}

\section*{Common Traps --- Worksheet \#1}
\begin{enumerate}[leftmargin=2em,itemsep=4pt]
  \item \textbf{Sign of the implicit gradient.} $6x-3y+15=0$ has gradient $2$, not $-2$ and not $3$ --- always solve for $y$ explicitly, or read $-a/b$ with both signs. \seealso{A1, A7, B1}
  \item \textbf{Distance: absolute value $\&$ square root together.} $\frac{|...+c|}{\sqrt{...}}$ --- dropping the absolute value floors the answer at zero and loses the sign of the offset. \seealso{A6, B2, C7}
  \item \textbf{Perpendicular does not mean ``negated slope''.} $m=2$ pairs with $-\frac12$, not $-2$; only the flip-invert passes $m_1m_2=-1$. \seealso{C5, C8, B1}
  \item \textbf{Shoelace: half the absolute difference, vertices in order.} Out-of-order vertices produce self-crossing regions whose signed sum misleads --- keep cyclic order and take half of the magnitude. \seealso{B4, B10, D5}
  \item \textbf{Collinearity: compare the right gradient pairs.} Three points are collinear iff {\em every} pair gives the same slope --- one matching pair is never enough. \seealso{A9, C6}
\end{enumerate}

\SpeedClosing{``Shoelace before altitude: area is a one-line cross product, not a construction.''}

\end{document}